\documentclass[11pt,a4paper]{article}
\usepackage{fontspec}
\setmainfont{NotoSerifCJKkr-Regular.otf}[Path=fonts/,BoldFont=NotoSerifCJKkr-Bold.otf,ItalicFont=NotoSerifCJKkr-Regular.otf]
\setsansfont{NotoSerifCJKkr-Regular.otf}[Path=fonts/,BoldFont=NotoSerifCJKkr-Bold.otf]
\setmonofont{DejaVu Sans Mono}[Scale=0.79]
\XeTeXlinebreaklocale "ko"
\XeTeXlinebreakskip=0pt plus 1pt
\usepackage[margin=23mm,top=23mm,bottom=23mm,headheight=15pt]{geometry}
\usepackage{amsmath,amssymb,amsthm,mathtools,mathrsfs}
\usepackage{booktabs,tabularx,longtable,array,graphicx,float}
\usepackage[dvipsnames]{xcolor}
\usepackage{enumitem,caption,fancyhdr,titlesec}
\usepackage{tikz}
\usetikzlibrary{arrows.meta,positioning,calc}
\usepackage{hyperref}
\providecommand{\tightlist}{\setlength{\itemsep}{0pt}\setlength{\parskip}{0pt}}
\newcommand{\passthrough}[1]{#1}
\hypersetup{unicode=true,colorlinks=true,linkcolor=MidnightBlue,urlcolor=MidnightBlue,citecolor=MidnightBlue,pdftitle={WU088 HH PHYS18: HH와 birth의 혼합 응답과 유한 상호작용},pdfauthor={WU088 HH 연구}}
\urlstyle{same}
\setlength{\emergencystretch}{3em}
\setlength{\parskip}{0.34em}
\setlength{\parindent}{1em}
\widowpenalty=10000\clubpenalty=10000
\linespread{1.10}
\setlist{itemsep=0.18em,topsep=0.35em,leftmargin=1.6em}
\captionsetup{font=small,labelfont=bf}
\renewcommand{\figurename}{그림}
\renewcommand{\tablename}{표}
\renewcommand{\contentsname}{목차}
\renewcommand{\refname}{참고 문헌과 원자료}
\renewcommand{\proofname}{증명}
\titleformat{\section}{\large\sffamily\bfseries\color{MidnightBlue}}{\thesection}{0.7em}{}
\titleformat{\subsection}{\sffamily\bfseries}{\thesubsection}{0.7em}{}
\pagestyle{fancy}\fancyhf{}
\fancyhead[L]{\small\sffamily WU088\_HH / PHYS18 이론}
\fancyhead[R]{\small\sffamily 2026-10-13 KST}
\fancyfoot[C]{\thepage}
\newcommand{\dd}{\mathrm d}
\newcommand{\diag}{\operatorname{diag}}
\newcommand{\ex}{\operatorname{exact64}}
\newcommand{\BB}{\mathcal B}
\newcommand{\HF}{\mathcal H_F}
\newcommand{\R}{\mathsf R}
\newcommand{\RF}{\mathsf R_F}
\newcommand{\ew}{e_w}
\newcommand{\exx}{e_x}
\newcommand{\nS}[1]{\left\lVert#1\right\rVert_S}
\newcommand{\nI}[1]{\left\lVert#1\right\rVert_\infty}
\newcommand{\magv}{\operatorname{mag}}
\newtheorem{theorem}{정리}
\newtheorem{lemma}{보조정리}
\newtheorem{proposition}{명제}
\theoremstyle{definition}\newtheorem{definition}{정의}
